%% kapa: exam=1 round=31 year=2020 period=1 form=A subject=civil_law subject_ko=민법 seq=13 number=13
%% source: Q-Net 감정평가사 기출문제 (cst00309) artlSeq=5207219 file=extracted/1008/5207219/A형/A형/2020년도 제31회 감정평가사 1차 시험_1교시_A형.pdf
%% conversion: kordoc 4.24.51 -> tools/gen1.py (auto); review=proofread-sonnet; flags=none
\begin{question}{13}
甲은 乙의 기망으로 그 소유의 X토지를 丙에게 팔았고, 丙은 그의 채권자 丁에게 X토지에 근저당권을 설정하였다. 甲은 기망행위를 이유로 매매계약을 취소하려고 한다. 이에 관한 설명으로 옳지 않은 것은? (다툼이 있으면 판례에 따름)
\begin{choices}
  \item 甲은 丙이 그의 잘못없이 기망사실을 몰랐을 때에만 매매계약을 취소할 수 있다.
  \item 丙의 악의 또는 과실은 甲이 증명하여야 한다.
  \item 甲은 매매계약을 취소하지 않고 乙에게 불법행위책임을 물을 수 있다.
  \item 丁의 선의는 추정된다.
  \item 매매계약을 취소한 甲은, 丁이 선의이지만 과실이 있으면 근저당권설정등기의 말소를 청구할 수 있다.
\end{choices}
\end{question}
\answer{1,5}
\answernote{복수 정답(정답표 원문: 1, 5)}
